\renewcommand{\headline}{Πέντε$\cdot$παιχνίδι - Eλληνικά}
\renewcommand{\tocent}{Eλληνικά}
\renewcommand{\tocline}{Κανόνες στα ελληνικά}
\renewcommand{\translator}{Athanasios Kostopoulos}
\renewcommand{\general}{
\lettrine[lines=3]{E}{να} παιχνίδι που μπορεί να εξηγηθεί σε ένα λεπτό αλλά παραμένει
συναρπαστικό για χρόνια.

Μια παρτίδα δύο παικτών διαρκεί μόλις 20 λεπτά. Μια παρτίδα
τριών ή τεσσάρων παικτών μπορεί να διαρκέσει ως 90 λεπτά.

Το παιχνίδι δεν χρησιμοποιεί ζάρια και είναι κατάλληλο για όλες τις
ηλικίες από 5 ετών και πάνω. 
%Ο σχεδιαστης του παιχνιδιου ειναι ο Jan \noun{Suchanek.}

Στο κουτί θα βρείτε: 4$\times$5 έγχρωμες φιγούρες, 5 μαύρα και
5 γκρίζα τουβλάκια, καθώς και το ταμπλό παιχνιδιού.
}

\renewcommand{\choosext}{Επιλογή φιγούρων}
\renewcommand{\choosex}{
Κάθε παίκτης έχει πέντε φιγούρες. Υπάρχει μία ομάδα φιγούρων ανά παίκτη στο κουτί.
%Μια ομαδα εχει ασπρα μαλλια, μια ομαδα μαυρα, μια ξανθα και μια ειναι φαλακρη.

%Οι παικτες πρεπει να διαλεξουν την ομαδα με την οποια μοιαζουν περισσοτερο.

Ο κάθε παίκτης έχει μία μπλε, μία κόκκινη, μία λευκή, μία πράσινη
και μία κίτρινη φιγούρα.

Αρχίζουν στις πέντε γωνίες του ταμπλό, ανάλογα με το χρώμα τους.

Ο στόχος είναι να φτάσουν στις μεγάλες στάσεις στο κέντρο του ταμπλό.

}
\renewcommand{\setupt}{Στήσιμο}
\renewcommand{\setup}{
Τοποθετήστε τις φιγούρες σας στις μεγάλες γωνίες του δακτυλίου, ανάλογα με το χρώμα τους:
η λευκή φιγούρα στη λευκή γωνία, η μπλε φιγούρα στην μπλε γωνία, κ.ο.κ.

Τοποθετήστε τα μαύρα τουβλάκια στις πέντε διασταυρώσεις στο μέσο του ταμπλό.

Αφήστε τα γκρίζα τουβλάκια στη μέση του ταμπλό για μετέπειτα χρήση. 
}

\renewcommand{\objectivet}{Στόχος παιχνιδιού}
\renewcommand{\objective}{
Οι λευκές φιγούρες θέλουν να πάνε στη λευκή διασταύρωση στο κέντρο,
οι μπλε φιγούρες στην μπλε, κ.ο.κ. Ο προορισμός είναι πάντα η μεγάλη,
έγχρωμη στάση στο μεσαίο πεντάγωνο, απέναντι από το σημείο εκκίνησης.

Ο πρώτος παίκτης που θα τοποθετήσει \emph{τρία} κομμάτια στον προορισμό τους, κερδίζει.
}

\renewcommand{\rulest}{Κανόνες παιχνιδιού}
\renewcommand{\rules}{
Μετακινήστε οποιαδήποτε από τις φιγούρες σας στο άστρο ή στο δακτύλιο, σε κάθε
κατεύθυνση, όσο θέλετε.

Μπορείτε να στρίψετε σε κάθε ελεύθερη γωνία ή διασταύρωση χωρίς να σταματήσετε.

Ποτέ δεν μπορείτε να πηδήξετε\textemdash ούτε πάνω από τουβλάκια, ούτε πάνω από φιγούρες. 

\myskip

Παρόλα αυτά, μπορείτε να μετακινηθείτε \emph{σε} μια κατειλημμένη θέση:

\myskip

Αν καταλάβετε θέση με μαύρο τουβλάκι, τοποθετήστε το σε μια ελεύθερη θέση της επιλογής σας.

Αν μετακινηθείτε σε θέση με μια άλλη φιγούρα, ανταλλάξτε θέσεις.
 
\myskip

Μπορείτε να ανταλλάξετε τις θέσεις δύο από τις φιγούρες σας.

Αν μετακινηθείτε σε θέση με πολλαπλές φιγούρες, διαλέξτε μία για ανταλλαγή θέσης.
 
\myskip 

Μην κάνετε την ίδια κίνηση δύο φορές συνεχόμενα.   

\myskip 

Μια φιγούρα που φτάνει στον προορισμό της αφαιρείται. Τοποθετήστε την στο κέντρο του ταμπλό. Αφού γίνει αυτό,
παίρνετε ένα από τα γκρίζα τουβλάκια και το τοποθετείτε σε ελεύθερη θέση της επιλογής σας.

\myskip

Αν πάτε σε θέση με γκρίζο τουβλάκι, το αφαιρείτε από το ταμπλό.

\myskip

Εάν μεταφερθείτε στον προορισμό σας από άλλο παίκτη, πρέπει να βγείτε από τη στροφή σας.

\myskip

Ο τελευταίος γύρος παίζεται μέχρι το τέλος.
}
